% p.28 — الباب الخامس العشر في رمى يرمي (ناقص يائي)
\mspage{28}
\bab{\com{الباب الخامس العشر}{كذا في الأصل؛ ولعل الصواب: الباب الخامس عشر} في رمٰى يرمي (ناقص يائي)}
% Layout: six tall columns (right-to-left), each a stack of boxed blocks of unequal height.
% Column 1 is the summary column: the باب name (ضَرَبَ) in its own box, then the ten
% underlying (pre-iʿlāl) forms رَمَيَ … لَا تَرْمِيْ, then the ten actual forms رَمٰى … لَا تَرْمِ.
% In columns 2–6 every sub-heading box holds the heading and, beneath it, the underlying form
% (e.g. «ماضي معروف / رَمَيَ»); the regular paradigm follows outside the box.
% Row alignment is approximate (all columns run to 21 lines here); \cline marks the printed box
% boundaries inside each column. Forms ending in alif maqṣūra carry a small (dagger) alif
% in the lithograph (رَمٰى، يُرْمٰى); it is transcribed as U+0670.
\begin{sarf}{|c|c|c|c|c|c|}
ضَرَبَ & \thead{ماضي معروف} & تَرْمِيْ & رُمِيَ & يُرْمَيْنَ & اِرْمُوْا\\\cline{1-1}
رَمَيَ & رَمَيَ & تَرْمِيَانِ & رُمِيَا & تُرْمٰى & اِرْمِيْ\\\cline{2-2}
يَرْمِيُ & رَمٰى & يَرْمِيْنَ & رُمُوْا & تُرْمَيَانِ & اِرْمِيَا\\
رَمْيًا & رَمَيَا & تَرْمِيْ & رُمِيَتْ & تُرْمَوْنَ & اِرْمِيْنَ\\\cline{6-6}
رَامِيٌ & رَمَوْا & تَرْمِيَانِ & رُمِيَتَا & تُرْمَيْنَ & \thead{نهي حاضر}\\
رُمِيَ & رَمَتْ & تَرْمُوْنَ & رُمِيْنَ & تُرْمَيَانِ & لَا تَرْمِيْ\\\cline{6-6}
يُرْمَيُ & رَمَتَا & تَرْمِيْنَ & رُمِيْتَ & تُرْمَيْنَ & لَا تَرْمِ\\
رَمْيًا & رَمَيْنَ & تَرْمِيَانِ & رُمِيْتُمَا & أُرْمٰى & لَا تَرْمِيَا\\
مَرْمُوْيٌ & رَمَيْتَ & تَرْمِيْنَ & رُمِيْتُمْ & نُرْمٰى & لَا تَرْمُوْا\\\cline{5-5}
اِرْمِيْ & رَمَيْتُمَا & أَرْمِيْ & رُمِيْتِ & \thead{اسم مفعول} & لَا تَرْمِيْ\\
لَا تَرْمِيْ & رَمَيْتُمْ & نَرْمِيْ & رُمِيْتُمَا & مَرْمُوْيٌ & لَا تَرْمِيَا\\\cline{1-1}\cline{3-3}\cline{5-5}
رَمٰى & رَمَيْتِ & \thead{اسم فاعل} & رُمِيْتُنَّ & مَرْمِيٌّ & لَا تَرْمِيْنَ\\\cline{6-6}
يَرْمِيْ & رَمَيْتُمَا & رَامِيٌ & رُمِيْتُ & مَرْمِيَّانِ & \thead{مضارع بلم}\\\cline{3-3}
رَمْيًا & رَمَيْتُنَّ & رَامٍ & رُمِيْنَا & مَرْمِيُّوْنَ & لَمْ يَرْمِيْ\\\cline{4-4}\cline{6-6}
رَامٍ & رَمَيْتُ & رَامِيَانِ & \thead{مضارع مجهول} & مَرْمِيَّةٌ & لَمْ يَرْمِ\\
رُمِيَ & رَمَيْنَا & رَامُوْنَ & يُرْمَيُ & مَرْمِيَّتَانِ & لَمْ يَرْمِيَا\\\cline{2-2}\cline{4-4}
يُرْمٰى & \thead{مضارع معروف} & رَامِيَةٌ & يُرْمٰى & مَرْمِيَّاتٌ & لَمْ يَرْمُوْا\\\cline{5-5}
رَمْيًا & يَرْمِيُ & رَامِيَتَانِ & يُرْمَيَانِ & \thead{أمر حاضر} & لَمْ تَرْمِ\\\cline{2-2}
مَرْمِيٌّ & يَرْمِيْ & رَامِيَاتٌ & يُرْمَوْنَ & اِرْمِيْ & لَمْ تَرْمِيَا\\\cline{3-3}\cline{5-5}
اِرْمِ & يَرْمِيَانِ & \thead{ماضي مجهول} & تُرْمٰى & اِرْمِ & لَمْ يَرْمِيْنَ\\
لَا تَرْمِ & يَرْمُوْنَ & رُمِيَ & تُرْمَيَانِ & اِرْمِيَا & لَمْ تَرْمِ\\
\end{sarf}
